\chapter{线性空间}
\orig{24--41}

\section{线性空间}

\subsection{定义及公理的独立性}

\textbf{1. 定义。} 设 $V$ 是一个非空集合，$F$ 是一个数域，$V$ 上有一个加法，$F$ 与 $V$ 之间有一个数乘。若满足以下条件，就称 $V$ 为 $F$ 上的一个\textbf{线性空间}。

加法满足
\begin{enumerate}[label=\arabic*),leftmargin=2.6em,itemsep=0.35em]
\item $\alpha+\beta=\beta+\alpha$；
\item $(\alpha+\beta)+\gamma=\alpha+(\beta+\gamma)$；
\item 在 $V$ 中有一个元素 $0$，对任意 $\alpha\in V$，有 $\alpha+0=\alpha$；
\item 对任意 $\alpha\in V$，存在 $\beta\in V$，使 $\alpha+\beta=0$。
\end{enumerate}
数乘满足
\begin{enumerate}[label=\arabic*),leftmargin=2.6em,itemsep=0.35em,start=5]
\item $1\alpha=\alpha$；
\item $k(l\alpha)=(kl)\alpha$。
\end{enumerate}
数乘与加法满足
\begin{enumerate}[label=\arabic*),leftmargin=2.6em,itemsep=0.35em,start=7]
\item $(k+l)\alpha=k\alpha+l\alpha$；
\item $k(\alpha+\beta)=k\alpha+k\beta$。
\end{enumerate}
其中 $k,l\in F$，$\alpha,\beta,\gamma\in V$。

根据这个定义，条件 1) 可由其他七条推出。只要分别计算
\[
(1+1)(\alpha+\beta)
 =(1+1)\alpha+(1+1)\beta
 =\alpha+(\alpha+\beta)+\beta
\]
以及
\[
(1+1)(\alpha+\beta)
 =1(\alpha+\beta)+1(\alpha+\beta)
 =\alpha+(\beta+\alpha)+\beta,
\]
即可得到
\[
\alpha+\beta=\beta+\alpha.
\]

但是其余七条 2)--8) 是相互独立的。为此，也可以采用下面与上述定义等价的公理组：加法满足
\[
\text{A) }\alpha+\beta=\beta+\alpha,
\qquad
\text{B) }(\alpha+\beta)+\gamma=\alpha+(\beta+\gamma),
\]
以及
\[
\text{C) 对任意 }\alpha,\beta\in V,
\quad \alpha+x=\beta\text{ 在 }V\text{ 中可解};
\]
数乘满足
\[
\text{D) }1\alpha=\alpha,
\qquad
\text{E) }k(l\alpha)=(kl)\alpha,
\]
且
\[
\text{F) }(k+l)\alpha=k\alpha+l\alpha,
\qquad
\text{G) }k(\alpha+\beta)=k\alpha+k\beta.
\]

下面逐条给出 A)--G 独立的例子。

\textbf{例 1（A 独立）。} $F$ 为数域，$V$ 为所有正有理数。规定
\[
\alpha+\beta=\beta,
\qquad
k\alpha=\alpha.
\]
容易验证 B)--G 均成立，但 A 不成立，例如
\[
2+3=3,
\qquad
3+2=2.
\]

\textbf{例 2（B 独立）。} 令 $V=\{0,\alpha,\beta\}$，加法表为
\[
\begin{array}{c|ccc}
+&0&\alpha&\beta\\\hline
0&0&\beta&\alpha\\
\alpha&\beta&\alpha&0\\
\beta&\alpha&0&\beta
\end{array}
\]
并对任意标量 $k\in F$、任意 $v\in V$ 规定 $kv=v$。这个运算交换、幂等，而且每一行都是 $V$ 的一个排列，所以 A 与 C 成立；D--G 也直接成立。但
\[
(0+\alpha)+\alpha=\beta+\alpha=0,
\qquad
0+(\alpha+\alpha)=0+\alpha=\beta,
\]
故 B 不成立。\jiaokan{第 26 页所列四元素表中，$\alpha$ 行只出现 $\alpha$ 与 $0$，因此方程 $\alpha+x=\beta$ 无解，条件 C 也不成立，不能据此证明 B 的独立性。这里改用满足 A、C、D--G 而只破坏 B 的三元素反例。}

\textbf{例 3（C 独立）。} $F$ 为数域，$V=\{0,\alpha,\beta\}$，加法表为
\[
\begin{array}{c|ccc}
+&0&\alpha&\beta\\\hline
0&0&\alpha&\alpha\\
\alpha&\alpha&\alpha&\alpha\\
\beta&\alpha&\alpha&\beta
\end{array}
\]
并规定 $k\alpha=\alpha$。则方程
\[
\alpha+x=\beta
\]
在 $V$ 中无解，而其余条件可验证成立。

\textbf{例 4（D 独立）。} 取 $F=\mathbb R$，$V=\mathbb R$，使用普通加法，并规定
\[
k\alpha=0.
\]
则 $1\alpha=\alpha$ 一般不成立，而其余各条成立。

\textbf{例 5（E 独立）。} 取
\[
F=\mathbb Q(\sqrt2)=\{a+b\sqrt2\mid a,b\in\mathbb Q\},
\qquad
V=\mathbb Q.
\]
使用普通加法，并规定
\[
(a+b\sqrt2)\alpha=(a+b)\alpha.
\]
由于
\[
(\sqrt2\sqrt2)\cdot1=2,
\qquad
\sqrt2(\sqrt2\cdot1)=1,
\]
故 E 不成立，而其余各条成立。

\textbf{例 6（F 独立）。} 取 $F=V=\mathbb R$，使用普通加法，并规定
\[
k\alpha=\alpha.
\]
显然一般有
\[
(k+l)\alpha\ne k\alpha+l\alpha,
\]
而其余各条成立。

\textbf{例 7（G 独立）。} 取 $F=\mathbb C$，$V=\mathbb C^2$，向量加法为通常加法。对每一条复直线 $L\subset V$，规定一个域自同构 $\sigma_L$：令
\[
L_0=\mathbb C(1,1),
\qquad
\sigma_{L_0}(z)=\bar z,
\]
而其余复直线上的 $\sigma_L$ 均取恒等自同构。对 $v\ne0$，若 $v\in L$，定义
\[
k\cdot v=\sigma_L(k)v,
\qquad  k\cdot0=0.
\]
在每一条直线上，D、E、F 均由 $\sigma_L$ 是域自同构而成立；A、B、C 来自 $V$ 的通常加法。但是取
\[
u=(1,0),\qquad v=(0,1),\qquad k=i,
\]
有
\[
i\cdot(u+v)=-i(1,1),
\qquad
 i\cdot u+i\cdot v=i(1,1),
\]
故 G 不成立。\jiaokan{第 27--28 页的例 7 规定“向量加法”为实数乘法、数乘仍为通常乘法，并声称只有 G 失败；但 F 也失败。例如取 $k=l=1,\alpha=3$，则 $(k+l)\alpha=6$，而 $k\alpha+_V l\alpha=3\cdot3=9$。为保持“各公理相互独立”的论证目的，这里以一个确实只破坏 G 的例子替换。}

\subsection{加群未必能作成域上的线性空间}

是不是每一个加群都能作成某个域上的线性空间？答案是否定的。一个加群能否作成域上的线性空间，数乘起着重要作用；也可能无论怎样规定数乘都不能作成线性空间。

例如令 $V$ 为整数加群，$\mathbb Z_p$ 为以素数 $p$ 为模的剩余类域。若 $V$ 能作成 $\mathbb Z_p$ 上的线性空间，则因 $\operatorname{char}\mathbb Z_p=p$，有
\[
\underbrace{1+1+\cdots+1}_{p\text{ 个}}=0.
\]
取 $2\in V$，由线性空间公理应有
\[
1\cdot2=2,
\qquad
0\cdot2=0.
\]
于是
\[
0=0\cdot2
 =\Bigl(\underbrace{1+1+\cdots+1}_{p\text{ 个}}\Bigr)\cdot2
 =\underbrace{2+2+\cdots+2}_{p\text{ 个}}
 =2p,
\]
但在整数加群中 $2p\ne0$，矛盾。因此无论怎样规定数乘，整数加群都不能作成 $\mathbb Z_p$ 上的线性空间。

\section{线性相关性}

\textbf{1. 定义。} 设 $\alpha_1,\ldots,\alpha_r$ 是线性空间 $V$ 的 $r$ 个向量。若在 $F$ 中存在不全为零的数 $k_1,\ldots,k_r$，使
\[
k_1\alpha_1+k_2\alpha_2+\cdots+k_r\alpha_r=0,
\]
则称 $\alpha_1,\ldots,\alpha_r$ \textbf{线性相关}。

\textbf{2.} $\beta$ 不能由 $\alpha_1,\ldots,\alpha_r$ 线性表示，并不保证 $\beta,\alpha_1,\ldots,\alpha_r$ 线性无关。

例如
\[
\beta=(1,1,1),\qquad \alpha_1=(1,0,0),\qquad \alpha_2=(2,0,0).
\]
明显 $\beta$ 不能由 $\alpha_1,\alpha_2$ 线性表示，但 $\beta,\alpha_1,\alpha_2$ 仍线性相关。这个例子也说明：一组向量线性相关，并不意味着其中每一个向量都能由其余向量线性表示。

\textbf{3.} 若 $\alpha_1,\ldots,\alpha_r$ 线性无关，则其中任意两个不同向量都线性无关；反之不真。取
\[
\alpha_1=(1,1,1),\quad \alpha_2=(0,1,0),\quad \alpha_3=(1,0,1).
\]
任意两个向量都线性无关，但
\[
\alpha_1=\alpha_2+\alpha_3,
\]
故三者线性相关。

\textbf{4.} 设 $\alpha_1,\ldots,\alpha_m$ 线性无关，考察
\[
\alpha_1+\alpha_2,\ \alpha_2+\alpha_3,\ \ldots,\ \alpha_{m-1}+\alpha_m,\ \alpha_m+\alpha_1.
\]
当 $m$ 为奇数且 $\operatorname{char}F\ne2$ 时，这 $m$ 个向量线性无关。事实上若
\[
\sum_{j=1}^{m-1}k_j(\alpha_j+\alpha_{j+1})+k_m(\alpha_m+\alpha_1)=0,
\]
由 $\alpha_1,\ldots,\alpha_m$ 线性无关得到
\[
k_1+k_m=0,
\quad k_1+k_2=0,
\quad \ldots,
\quad k_{m-1}+k_m=0.
\]
其系数行列式为
\[
D=1-(-1)^m=2
\]
（$m$ 为奇数），故在 $\operatorname{char}F\ne2$ 时 $D\ne0$，只有零解。

当 $m$ 为偶数时，上述向量一定线性相关，因为
\[
(\alpha_1+\alpha_2)-(\alpha_2+\alpha_3)+\cdots-(\alpha_m+\alpha_1)=0.
\]
\jiaokan{由 $D=2\ne0$ 推出线性无关需要排除特征 $2$。若把 $F$ 理解为一般域，则特征 $2$ 时结论不成立；旧式“数域”通常默认特征 $0$，在该语境下结论成立。}

\section{子空间}

\textbf{1. 定义。}

1) 数域 $F$ 上线性空间 $V$ 的一个非空子集 $S$，若对 $V$ 的加法与数乘也构成线性空间，则称 $S$ 为 $V$ 的一个\textbf{子空间}。

2) 若 $S_1,S_2$ 是 $V$ 的子空间，定义它们的和为
\[
S_1+S_2=\{\alpha_1+\alpha_2\mid \alpha_1\in S_1,\ \alpha_2\in S_2\}.
\]

3) 若 $S_1+S_2$ 中每个向量的表示
\[
\alpha=\alpha_1+\alpha_2,
\qquad \alpha_i\in S_i\ (i=1,2)
\]
都是唯一的，则称 $S_1+S_2$ 为\textbf{直和}。

\textbf{2.} 子空间的直和当然是子空间的和，但子空间的和未必是直和。

例如取 $F=\mathbb R$，
\[
V=\{(a_1,a_2,a_3)\mid a_i\in\mathbb R\},
\]
\[
S_1=\{(a_1,a_2,0)\mid a_i\in\mathbb R\},
\qquad
S_2=\{(0,b_2,b_3)\mid b_i\in\mathbb R\}.
\]
显然 $V=S_1+S_2$。对 $(x,y,z)\in V$，有
\[
(x,y,z)=(x,y,0)+(0,0,z)
       =(x,0,0)+(0,y,z).
\]
只要 $y\ne0$，这就是两种不同表示，所以 $S_1+S_2$ 不是直和。

\textbf{3.} 设 $S_1,\ldots,S_n\ (n\ge4)$ 是 $V$ 的子空间，并满足
\[
S_i\cap(S_j+S_k)=0\qquad(i\ne j,k).
\]
也不能推出 $S_1+\cdots+S_n$ 一定是直和。

例如在 $F=\mathbb R$、$V=\mathbb R^3$ 中令
\[
S_1=L[(1,0,0)],\quad
S_2=L[(0,1,0)],\quad
S_3=L[(0,0,1)],\quad
S_4=L[(-1,-1,-1)].
\]
容易验证 $S_i\cap(S_j+S_k)=0$（$i\ne j,k$），但
\[
(1,0,0)+(0,1,0)+(0,0,1)+(-1,-1,-1)=0,
\]
故 $S_1+S_2+S_3+S_4$ 不是直和。

\section{环上的线性空间}

在数域 $F$ 上的线性空间中，维数是唯一的，而且子空间的维数不超过原线性空间的维数。因此有限维空间的子空间仍为有限维。把数域换成环以后，这些结论未必成立。下面给出若干例子。

\subsection{定义}

以下的环 $R$ 均指有单位元的环。

1) \textbf{环上的线性空间。} 一个加群 $V$ 若配有 $R$ 到 $V$ 的数乘，并满足
\[
1\alpha=\alpha,
\qquad
k(l\alpha)=(kl)\alpha,
\]
\[
(k+l)\alpha=k\alpha+l\alpha,
\qquad
k(\alpha+\beta)=k\alpha+k\beta,
\]
则称 $V$ 为环 $R$ 上的一个左线性空间。这里 $k,l\in R$，$\alpha,\beta\in V$。

2) \textbf{循环空间。} 若 $V$ 恰好包含所有 $k\alpha\ (k\in R)$，记
\[
R\alpha=\{k\alpha\mid k\in R\},
\]
则称其为由 $\alpha$ 生成的循环空间，$\alpha$ 称为一个生成元。

3) \textbf{有限空间。} 若 $V$ 是有限多个循环子空间
\[
R\alpha_1,\ldots,R\alpha_n
\]
的和，则称 $V$ 为 $R$ 上的一个有限空间。

4) \textbf{有基底空间。} 若 $V$ 是 $R\alpha_1,\ldots,R\alpha_n$ 的直和，则称 $\alpha_1,\ldots,\alpha_n$ 为 $V$ 的一个基底，$V$ 称为有基底空间。

5) \textbf{自由空间。} 若上述直和还满足 $0$ 是每个 $\alpha_i$ 唯一的零化子，即
\[
k\alpha_i=0\Longrightarrow k=0,
\]
则称 $V$ 为自由空间，$\alpha_1,\ldots,\alpha_n$ 为自由基底，并把 $n$ 称为这个自由基底的长度。\jiaokan{本节沿用“有限空间”“有基底空间”“自由空间”等旧式术语；它们分别更接近现代模论中的“有限生成模”“若干循环子模的内直和”与“自由模”。现代教材通常把“基底”保留给自由模，这里的术语与今天的通行用法并不完全一致。}

\textbf{2.} 整数加群不能作成 $\mathbb Z_p$ 上的线性空间，见本章第 1 节的例子。

\textbf{3.} 环上线性空间中，子空间和中的元素表示法未必唯一。

例如 $R=V=\mathbb Z$，把 $V$ 看作 $R$ 上的线性空间。令
\[
S_1=(4),\qquad S_2=(6).
\]
则 $10\in S_1+S_2$，且
\[
10=4+6=16+(-6),
\]
有两种不同的表示。

\textbf{4.} 有限空间未必是有基底空间。

例如 $\mathbb Z[x]$ 是整数环 $\mathbb Z$ 上的多项式环，理想
\[
(2,x)
\]
是 $\mathbb Z[x]$ 上由 $2,x$ 生成的有限空间，但不是有基底空间。

若 $(2,x)$ 有基底 $f_1(x),\ldots,f_n(x)$。若其中有两个非零元 $f_1,f_2$，则
\[
0\ne f_1f_2\in\mathbb Z[x]f_1,
\qquad
0\ne -f_1f_2\in\mathbb Z[x]f_2,
\]
从而
\[
f_2f_1+(-f_1)f_2=0
\]
给出非平凡关系，和不可能是直和。于是基底中至多有一个非零元，这又意味着 $(2,x)$ 是主理想；但 $(2,x)$ 不是主理想，矛盾。

\textbf{5.} 有基底空间未必是自由空间。

取 $V=\mathbb Z_p$（$p$ 为素数），$R=\mathbb Z$，规定数乘为整数倍。则 $V$ 是 $R$ 上的有基底空间：
\[
V=R\cdot[1],
\]
而 $[1]$ 是一个基底。但 $[1]$ 的零化子理想为 $(p)$，并非只有 $0$；因此 $V$ 不是自由空间。

\textbf{6.} 自由空间可以有长度不同的自由基底。

设 $F$ 为域，$R$ 是所有 $F$ 上“每行、每列只有有限多个非零元”的无限矩阵所成的环，其单位元为无限单位矩阵 $E$。令 $V=R$，按左乘看作 $R$ 上的自由空间。把矩阵的奇数行、偶数行分别抽取并重新编号，定义 $K_1,K_2$；相应地，把奇数列、偶数列重新嵌回原列位置定义 $K_1',K_2'$。具体地，$K_1$ 的第 $r$ 行在第 $\lceil r/2\rceil$ 列为 $1$（$r$ 为奇数时），偶数行为 $0$；$K_2$ 则在偶数行做同样的抽取。相应的 $K_1',K_2'$ 满足
\[
(K_1,K_2)
\binom{K_1'}{K_2'}=E,
\]
以及
\[
\binom{K_1'}{K_2'}(K_1,K_2)
=\begin{pmatrix}E&O\\O&E\end{pmatrix}.
\]
因此 $E$ 是 $V$ 的一个自由基底，长度为 $1$；另一方面 $K_1',K_2'$ 也构成一个自由基底，长度为 $2$。这给出自由基底长度不唯一的例子。

\textbf{7.} 自由空间与其子空间的自由基底可以有相同长度。

取 $R=\mathbb Z$，$V$ 是 $R$ 上自由空间，$\alpha_1,\ldots,\alpha_n$ 为一个自由基底。令 $S$ 为 $2\alpha_1,\ldots,2\alpha_n$ 生成的子空间，则 $2\alpha_1,\ldots,2\alpha_n$ 仍是 $S$ 的自由基底，所以 $S$ 的长度仍为 $n$。

\textbf{8.} 环 $R$ 上的有限空间可以有不是有限空间的子空间。

设 $F$ 为域，
\[
R=F[x_1]\cup F[x_1,x_2]\cup F[x_1,x_2,x_3]\cup\cdots,
\]
即有限多个不定元 $x_1,x_2,\ldots$ 上的多项式环之并。令 $V=R$，则 $V$ 是 $R$ 上的有限空间，以 $1$ 为自由基底。令
\[
S=\{0\}\cup\{\text{所有不含常数项的多项式}\}.
\]
显然 $S$ 是 $V$ 的子空间，但不是 $R$ 上的有限空间。否则若 $S$ 由
\[
\alpha_i=f_i(x_{i1},\ldots,x_{i r_i})\qquad(i=1,\ldots,n)
\]
有限生成，取一个没有出现在任何这些多项式中的不定元 $x_k$。若
\[
x_k=g_1\alpha_1+\cdots+g_n\alpha_n,
\qquad g_i\in R,
\]
令 $X$ 为所有出现在 $f_1,\ldots,f_n$ 中的不定元的有限集合，并取 $x_k\notin X$。考虑把 $X$ 中所有不定元都送到 $0$、而保持 $x_k$ 不变的代入同态。由于每个 $f_i$ 都无常数项且只含 $X$ 中的不定元，代入后 $f_i$ 全部变为 $0$；于是等式右端变为 $0$，左端仍为 $x_k$，矛盾。

\textbf{9.} 自由空间的子空间可以不是有限空间。

取一个有非有限生成左理想 $A$ 的环 $R$。令
\[
V=R\oplus R,
\qquad
S=\{(a,0)\mid a\in A\}.
\]
$V$ 是 $R$ 上的自由空间，$S$ 是其子空间。若 $S$ 有有限生成元
\[
\alpha_i=(a_i,0),\qquad i=1,\ldots,n,
\]
则任意 $(x,0)\in S$ 都可写成
\[
(x,0)=l_1\alpha_1+\cdots+l_n\alpha_n
=(l_1a_1+\cdots+l_na_n,0),
\]
于是
\[
A=Ra_1+\cdots+Ra_n,
\]
与 $A$ 不是有限生成左理想矛盾。

\textbf{10.} 自由空间的子空间未必是自由空间。

取有零因子的环 $R=\mathbb Z_n$，其中 $n=2m$ 为偶数。令 $V=\mathbb Z_n$，则 $V$ 是 $R$ 上的自由空间。令
\[
S=\{[0],[2],\ldots,[n-2]\}.
\]
则 $S$ 是 $V$ 的子空间，由 $[2]$ 生成，但不是自由空间，因为
\[
[2][m]=[2m]=[n]=[0]
\]
而 $[m]\ne[0]$。例如在 $\mathbb Z_6$ 中，
\[
S=\{[0],[2],[4]\},
\qquad [2][3]=[0].
\]

\textbf{11.} 把环 $R$ 上的左线性空间 $V$ 中的数乘形式地改写为
\[
\alpha k=k\alpha
\]
并不总能得到一个右线性空间；当 $R$ 非交换时尤其如此。

例如取
\[
R=\left\{\begin{pmatrix}a&0\\b&c\end{pmatrix}\middle|a,b,c\in\mathbb Z_2\right\},
\]
这是一个有单位元的非交换环，共有 $8$ 个元素。令
\[
V=M_2(\mathbb Z_2),
\]
以矩阵加法为加法，以左乘为数乘，则 $V$ 是 $R$ 上的左线性空间。若硬性规定右作用 $\alpha k:=k\alpha$，则右模结合律要求
\[
\alpha(kl)=(\alpha k)l.
\]
取
\[
k=\begin{pmatrix}1&0\\0&0\end{pmatrix},
\qquad
l=\begin{pmatrix}1&0\\1&0\end{pmatrix},
\qquad
\alpha=\begin{pmatrix}1&0\\0&1\end{pmatrix},
\]
有
\[
\alpha(kl)=(kl)\alpha=\begin{pmatrix}1&0\\0&0\end{pmatrix},
\]
而
\[
(\alpha k)l=l(k\alpha)=\begin{pmatrix}1&0\\1&0\end{pmatrix}.
\]
二者不等，因此不能这样把该左线性空间变成右线性空间。
